\section*{Hoofdstuk \ref{ch:b1:electrophilic-additions} --- Alkenen en alkynen: elektrofiele addities}

\begin{solution}{exo:b1:electrophilic-additions:1}
2-Chloorpropaan; 2-broom-2-methylpropaan; 1-broom-1-methylcyclohexaan;
2-chloorbut-2-een.
\end{solution}

\begin{solution}{exo:b1:electrophilic-additions:2}
1,2-Dibroomethaan; \textit{trans}-1,2-dibroomcyclohexaan;
1,2-dibroompropaan.
\end{solution}

\begin{solution}{exo:b1:electrophilic-additions:3}
Propaan-2-ol; 2-methylpropaan-2-ol; 1-methylcyclohexaan-1-ol (telkens
Markovnikov).
\end{solution}

\begin{solution}{exo:b1:electrophilic-additions:4}
Propyn (\omterm{def:b1:electrophilic-additions:acetylide}{terminaal alkyn}): \ce{CH3C#CH + NH2- -> CH3C#C- + NH3}. Ethanol:
\ce{CH3CH2OH + NH2- -> CH3CH2O- + NH3}. But-2-yn en propeen hebben geen
voldoende zuur waterstofatoom.
\end{solution}

\begin{solution}{exo:b1:electrophilic-additions:5}
Het $\pi$-paar neemt \ce{H+} op op het \ce{CH2}-koolstofatoom: tertiair
kation \ce{(CH3)3C+}; een \omterm{def:b1:lewis-resonance-vsepr:lewis}{vrij paar} van water bindt eraan:
\ce{(CH3)3COH2+}; een watermolecuul neemt een proton op: \ce{(CH3)3COH}.
In heet geconcentreerd zuur is water schaars en ontsnapt het vluchtige
alkeen: de omgekeerde reactie, de \omterm{def:b1:alcohol-activation:dehydration}{dehydratatie}, verloopt.
\end{solution}

\begin{solution}{exo:b1:electrophilic-additions:6}
Protonering van 2-methylpropeen geeft een tertiair kation, die van
propeen een secundair. De snelheidsbepalende protonering heeft een
\omterm{def:b1:elementary-steps:energy-profile}{overgangstoestand} die op het kation lijkt (Hammond): het stabielere
tertiaire kation wordt over een lagere barrière bereikt.
\end{solution}

\begin{solution}{exo:b1:electrophilic-additions:7}
Het bromoniumion overbrugt één zijde van de ring; bromide opent het vanaf
de andere zijde: de twee broomatomen eindigen \textit{trans}. Het
\textit{trans}-dibromide is \omterm{def:b1:stereochemistry-in-depth:stereocentre}{chiraal} (geen spiegelvlak), maar bromide valt
de twee koolstofatomen even vaak aan: de twee \omterm{def:b1:stereochemistry-in-depth:enantiomers}{enantiomeren} ontstaan in
gelijke hoeveelheden, en het product is racemisch, optisch inactief.
\end{solution}

\begin{solution}{exo:b1:electrophilic-additions:8}
Protonering van C1 geeft het secundaire kation op C2, naast het tertiaire
C3 dat een waterstofatoom draagt; dat waterstofatoom verschuift met zijn
paar naar C2 (een 1,2-hydrideverschuiving), wat een tertiair kation op C3
achterlaat, dat door \ce{Cl-} ingevangen wordt: 2-chloor-2-methylbutaan.
\end{solution}

\begin{solution}{exo:b1:electrophilic-additions:9}
Propyn + natriumamide: het propynide-ion; met 1-broompropaan (SN2):
\ce{CH3C#C-CH2CH2CH3}, hex-2-yn. Propynide met benzaldehyde, daarna
verdund zuur: \ce{C6H5CH(OH)C#CCH3}, 1-fenylbut-2-yn-1-ol.
\end{solution}

\begin{solution}{exo:b1:electrophilic-additions:10}
\cip{E}-but-2-een: bromonium aan één zijde, bromide vanaf de andere
zijde op C2 of op C3; beide aanvallen geven dezelfde verbinding, met
descriptoren \cip{2R,3S}: de \omterm{def:b1:stereochemistry-in-depth:enantiomers}{mesoverbinding}, één \omterm{def:b1:stereochemistry-in-depth:isomers}{stereo-isomeer}.
\cip{Z}-but-2-een: aanval op C2 geeft \cip{2S,3S} (of zijn spiegelbeeld
vanaf de andere zijde), aanval op C3 \cip{2R,3R}; gelijke hoeveelheden:
twee \omterm{def:b1:stereochemistry-in-depth:isomers}{stereo-isomeren}, een \omterm{def:b1:stereochemistry-in-depth:enantiomers}{racemisch mengsel}.
\end{solution}

\begin{solution}{exo:b1:electrophilic-additions:11}
In het asymmetrische bromoniumion draagt het meest gesubstitueerde
koolstofatoom een groter deel van de positieve lading (zijn C--Br-binding
is langer en zwakker, zoals bij een bijna tertiair of secundair kation).
Water, in grote overmaat, valt dat koolstofatoom aan vanaf de kant
tegenover de brug: 1-broompropaan-2-ol.
\end{solution}

\begin{solution}{exo:b1:electrophilic-additions:12}
$D = 1$: een alkeen. 2-Methylbut-1-een en 2-methylbut-2-een geven allebei
het tertiaire kation, en dus 2-broom-2-methylbutaan en
2-methylbutaan-2-ol. Hun protonen-NMR-spectra verschillen: twee
vinylprotonen (\ce{=CH2}, singuletten) bij het eerste, één (\ce{=CH-},
een kwartet) bij het tweede.
\end{solution}

\begin{solution}{pb:b1:electrophilic-additions:1}
\textbf{1.} Op elk koolstofatoom gaat \ce{CH3} vóór H: bij \cip{Z} staan
de twee methylgroepen aan dezelfde kant, bij \cip{E} aan tegengestelde
kanten.
\textbf{2.} \omterm{def:b1:stereochemistry-in-depth:isomers}{Stereo-isomeren} die geen spiegelbeelden zijn: \omterm{def:b1:stereochemistry-in-depth:enantiomers}{diastereomeren}.
\textbf{3.} Rotatie om C=C zou de $\pi$-binding verbreken, wat veel meer
energie kost dan botsingen bij kamertemperatuur leveren.
\textbf{4.} \cip{E}: zijn methylgroepen zitten elkaar niet in de weg.
\textbf{5.} \ce{C4H8 + Br2 -> C4H8Br2}.
\textbf{6.} Twee equivalente stereocentra: hoogstens vier, in
werkelijkheid drie (\cip{2R,3R}, \cip{2S,3S} en de mesovorm \cip{2R,3S}).
\textbf{7.} Het $\pi$-paar valt één Br-atoom van \ce{Br2} aan, het
Br--Br-paar vertrekt als \ce{Br-}, en een \omterm{def:b1:lewis-resonance-vsepr:lewis}{vrij paar} van het aangevallen
broomatoom bindt aan het tweede koolstofatoom: een bromoniumion met drie
ringatomen.
\textbf{8.} In het overbrugde ion heeft elk atoom een octet; een open
secundair \omterm{def:b1:electronic-effects:intermediates}{carbokation} zou een koolstofatoom met zes elektronen hebben.
\textbf{9.} Een \omterm{def:b1:lewis-resonance-vsepr:lewis}{vrij paar} van \ce{Br-} bindt aan één ringkoolstofatoom; de
C--Br-binding van de brug breekt, en haar paar gaat naar het
overbruggende broomatoom.
\textbf{10.} De brug bezet één zijde; zoals bij SN2 komt het \omterm{def:b1:electronic-effects:nucleophile}{nucleofiel}
binnen tegenover de binding die breekt.
\textbf{11.} Een \omterm{def:b1:electrophilic-additions:halonium}{anti-additie}.
\textbf{12.} Water zou met bromide om het bromoniumion concurreren en een
broomhydrine geven.
\textbf{13.} \cip{2R,3S} (of \cip{2S,3R}, dezelfde verbinding).
\textbf{14.} Dezelfde verbinding.
\textbf{15.} In een \omterm{def:b1:stereochemistry-in-depth:isomers}{conformatie} met de twee broomatomen eclips loopt er
een spiegelvlak tussen C2 en C3; de descriptoren zijn tegengesteld.
\textbf{16.} \cip{2R,3R} en \cip{2S,3S}.
\textbf{17.} Gelijke hoeveelheden: een \omterm{def:b1:stereochemistry-in-depth:enantiomers}{racemisch mengsel}.
\textbf{18.} Ja: de twee diastereomere alkenen geven verschillende
producten (meso tegenover racemisch).
\textbf{19.} Geen: de twee \omterm{def:b1:stereochemistry-in-depth:enantiomers}{enantiomeren} heffen elkaar op.
\textbf{20.} $5.60/56.0 = \qty{0.100}{mol}$ van elk.
\textbf{21.} $0.100 \times 0.85 \times 215.8 = \qty{18.3}{g}$ van elk.
\textbf{22.} \textbf{$[\alpha] = 0^\circ$ (\omterm{def:b1:stereochemistry-in-depth:enantiomers}{mesoverbinding}), \qty{18.3}{g}}.
\end{solution}
